\documentclass[10pt,a4paper]{amsart}
\usepackage[margin=19mm]{geometry}
\usepackage{amsmath,amssymb,mathtools}
\usepackage[scr=rsfs,cal=euler]{mathalfa}
\usepackage{xfrac}
\usepackage[unicode,hidelinks,bookmarks=false]{hyperref}
\newcommand{\ie}{i.e.\ }
\newcommand{\cf}{cf.\ }
\newcommand{\resp}{respectively\ }
\newcommand{\Subsection}{\subsection}
\providecommand{\nobreakdash}{\nobreak\hbox{-}}
\setlength{\emergencystretch}{2em}
\multlinegap0pt
\usepackage{amssymb}
\usepackage{float}
\newtheorem{theorem}{Theorem}
\title{Brik's sequence: a strange recursion}
\author{Jeffrey Shallit}
\date{Source: arXiv:2605.07542v1 (CC BY 4.0)}
\begin{document}
\maketitle

\noindent\emph{Source and licence.} Brik's sequence: a strange recursion. Version arXiv:2605.07542v1 (CC BY 4.0), distributed by arXiv under the Creative Commons Attribution 4.0 International licence, http://creativecommons.org/licenses/by/4.0/, as recorded in the arXiv licence metadata for this exact version and shown on the abstract page https://arxiv.org/abs/2605.07542v1. Full licence notice: \texttt{licenses/arxiv-2605.07542-CC-BY-4.0.txt}.\par\smallskip

\noindent\textit{Editorial scope.} Counted unit: the article's theorem theorem (source0.tex lines 278--281). The article's complete original proof of it is delivered with it (source0.tex lines 283--383, one \texttt{proof} environment of 101 lines). No mathematics is written, repaired or re-derived: the statement and the proof are the article's own lines, in the article's own notation, and no retained line is substituted or re-flowed.  No further source-local context is needed: every cross-reference of every retained line resolves inside the two delivered spans.  Nothing was omitted from any retained span, so the package carries no omission record.  No retained line contains a citation, so the wrapper carries no bibliography and no bibliography entry.  No retained line contains a figure environment, an image-inclusion command or drawing code, so no image is delivered and the package declares no asset dependency.  Editorial formatting only, in addition: an AMS \texttt{amsart} preamble that restates the article's own declarations and macros used by the retained lines, namely the environments \texttt{theorem} on separate counters, together with the standard packages those lines need.  The retained environments print in this document as Theorem 1 in order of appearance, whereas the article prints the same items with the numbering of its own full document.  The complete native source of the article as distributed in the arXiv e-print for this exact version is bundled unchanged as \texttt{sources/200089/source0.tex}.
\begin{theorem}
The limit $\alpha = \lim_{n \rightarrow \infty} a(n)/n$ exists, and
equals $2 - 2 \sum_{i \geq 1} b_i 2^{-i}$.
\end{theorem}

\begin{proof}
First we study the number $s(n)$ of $1$'s in prefixes of length
$\ell_n = |B_n|$.

Since $B_n = B_{n-1} B_{n-1}[n..\ell{n-1}]$, the number of $1$'s
the second factor is $s(n-1) - a(n-1)$.
Hence $s(n) = 2s(n-1) - a(n-1)$.  Dividing by $2^{n-1}$, we get
$$ {{s(n)}\over {2^{n-1}}} = {{s(n-1)} \over {2^{n-2}}} - 
{{a(n-1)} \over 2^{n-1}} .$$
Since $s(1) = 2$, iteration gives
\begin{equation}
 {{s(n)} \over {2^{n-1}}} = 2 - \sum_{1 \leq j < n} {{a(j)} \over {2^j}} .
\label{sn}
\end{equation}

Define $\alpha = 2 - \sum_{j \geq 1} {{a(j)} \over {2^j}} $.  This
is well-defined because $a(j) \leq j$ and
$\sum_{j \geq 1} j/2^j$ converges.   Then Eq.~\eqref{sn}
gives $\lim_{n \rightarrow \infty} {{s(n)}\over {2^{n-1}}} = \alpha$.
Since $\ell_n = 2^{n-1} + O(n)$, we also get
$\lim_{n \rightarrow \infty} {{s(n)}\over {\ell_n}} = \alpha$.

Next, we obtain an expression for $\alpha$.
Since $a(j) = \sum_{1 \leq i \leq j} b_i$,
we have
\begin{align*}
\sum_{j \geq 1} {{a(j)} \over {2^j}} &= \sum_{j\geq 1} 2^{-j} \sum_{1 \leq i \leq j} b_i \\
&= \sum_{i \geq 1} b_i \sum_{j \geq i} 2^{-j} \\
&= \sum_{i \geq 1} b_i 2^{1-i}  = 2 \sum_{i \geq 1} b_i 2^{-i}.
\end{align*}
Now if $\beta = \sum_{i \geq 1} b_i 2^{-i}$  then
\begin{equation}
\alpha = 2 - \sum_{j \geq 1} {{a(j)}\over {2^j}} = 2-2 \beta.
\label{alphabeta}
\end{equation}

We have now proven that the limiting density of 
$1$'s in $\bf b$ exists and equals
$\alpha$, if we only consider prefixes of length $\ell_i$.
It remains to consider arbitrary prefixes.

Define $E(N) = a(N) - \alpha N$ for $N \geq 1$.
From Eq.~\eqref{sn} we have
$$ {{s(n)} \over 2^{n-1}} = \alpha + \sum_{j \geq n} {{a(j)} \over
{2^j}} $$
and 
$$ \sum_{j \geq n} {{a(n)} \over {2^j}} = O(n 2^{-n}).$$
Therefore
$s(n) = \alpha 2^{n-1} + O(n)$.
Since $\ell_n = 2^{n-1} + n + 1$, it follows that
$$E(\ell_n) = s(n) - \alpha \ell_n = O(n).$$

Now consider an arbitrary integer $N \geq 1$.  
Write $N$ uniquely as $N = \ell_n + t$ where
$0 \leq t \leq \ell_{n+1} - \ell_n$.
Now
$\ell_{n+1} - \ell_n = \ell_n - n$, so 
$0 \leq t \leq \ell_n - n$.
Since $B_{n+1} = B_n B_n[n..\ell_n]$,
the length-$(\ell_n + t)$ prefix of $\bf b$
is $B_n B_n[n+1..n+t]$.  
Thus we get
$ a(N) = a(\ell_n + t) = s(n) + a(n+t)-a(n)$.
Subtracting $\alpha(\ell_n + t)$, we get
\begin{align}
E(\ell_n + t) &= s(n) - \alpha \ell_n + a(n+t) - a(n) - \alpha t  \nonumber \\
&= E(\ell_n) + E(n+t) - E(n)  \label{ee}.
\end{align}

Now define $\rho = \limsup_{N \rightarrow \infty} |E(N)|/N$.
From  the definition of $E$ we have $\rho < \infty$.

For $N = \ell_n + t$ we get from Eq.~\eqref{ee} that
$|E(N)| \leq |E(\ell_n)| + |E(n+t)| + |E(n)| $.
Dividing by $N = \ell_n + t$ we get
\begin{equation}
{{|E(N)|} \over N} \leq 
{{|E(\ell_n)|} \over N} + 
{{|E(n+t)|} \over {n+t}} \cdot {{n+t} \over {\ell_n + t}} +
{{|E(n)|} \over N}. \label{ineq}
\end{equation}
The first term tends to $0$ because
$|E(\ell_n)| = O(n)$.
The third term also tends to $0$
because $|E(n)| = O(n)$.

For the middle term,  note that 
$$ {{n+t} \over {\ell_n + t}} \leq {{\ell_n} \over {2\ell_n - n}} $$
because $0 \leq t \leq \ell_n - n$.
Now $(n+t)/(\ell_n + t)$ is an increasing function of $t$,
and its maximum on the interval $t \in [0, \ell_n - n]$
is attained at $t = \ell_n -n$, where it takes the value
$\ell_n/(2\ell_n - n)$.
Therefore this maximum is $1/2 + o(1)$.  

Taking $\limsup$ in \eqref{ineq}, we get $\rho \leq \rho/2$,
and hence $\rho = 0$.   Thus
$\lim_{N \rightarrow \infty} E(N)/N = 0$
and hence
$\lim_{N \rightarrow \infty} a(N)/N = \alpha$.
\end{proof}

\end{document}
